% !TEX root = ../../../main.tex
% Block card: preprocessing.window_pool -- "Window pool" (Part I, the blocks chapter).
\cardfield{What it does}
Starts a chain from window sets you saved in the Library, combines them into one pool, and
gives every window a role: train, validation, test or exam.

\cardfield{When to use it}
When you want to cluster or train on windows from many recordings, channels and window
lengths at once, as research question 1 does: pool the windows, cluster the training ones on
their shape (\blockref{preprocessing.trace_shape}, then \blockref{catalogue.shape_cluster}),
and test on time the clustering never saw. It is a \emph{source} block: it replaces the
span, so it can only be step 01. You make the window sets first, in the Library
(\apppage{library.window-sets}).

\begin{plainterms}[Roles before windows]
The pool decides the roles first, as stretches of time, and only then looks at the windows.
Think of a field split into plots before anyone plants: a plant belongs to the plot it
stands wholly inside. Each recording is cut into blocks of time; the last ones are the test,
the one before them validation, the rest training. Between two roles a strip of
no-man's-land is left empty, so no window of one role sits next to a window of another.
A window lying in that strip, or across a border, is dropped and counted. A whole mushroom
pack, or chosen channels, can be held out as the \emph{exam}: windows nothing is ever
trained or tuned on.
\end{plainterms}

\cardfield{What goes in, what comes out}
In go the saved window sets you tick (by default every saved unlabelled set). The signal of
the chain is ignored. Out comes one set of windows, each still knowing its recording,
channel, length (scale) and role. The pool is saved in the Library too; the same pool built
twice is reused, not saved again.

\cardfield{Settings that matter}
\emph{hold\_out\_pack} (default D, channels 13 to 16) is held out as the exam on every
recording. \emph{gap\_min} is the no-man's-land, in minutes; it must be at least the
longest window in the pool. \emph{per\_scale} draws an equal number of windows from each
window length (20,000 by default). \emph{pool} reopens a saved pool and ignores every other
setting.

\cardfield{Watch out for}
Raising \emph{gap\_min} or changing the blocks changes which windows survive; the counts of
windows dropped, and why, are on the block's page. The held-out recording (M4) is refused
outright, whatever you tick.
